\chapter*{Resumen}
\addcontentsline{toc}{chapter}{Resumen}
El presente Proyecto Integrador describe el diseño, la construcción y la validación experimental de un prototipo de granja urbana vertical inteligente, desarrollado en el marco de la línea de investigación del LIADE (Laboratorio de Investigación Aplicada y Desarrollo Electrónico) sobre agricultura urbana y periurbana. El objetivo del sistema es automatizar el cultivo en interiores mediante el monitoreo y control autónomo de las principales variables ambientales --temperatura, humedad del aire, humedad del sustrato e intensidad lumínica-- sin requerir conocimientos técnicos especializados por parte del usuario. Se diseñó y construyó una estructura física modular de varios niveles para el cultivo en sustrato tradicional, empleando rúcula y espinaca como cultivos modelo, seleccionados por su adaptabilidad agronómica a condiciones de interior. Se seleccionó una Raspberry Pi 4 como unidad de procesamiento central, encargada de recibir las lecturas de sensores digitales y analógicos y de accionar los actuadores mediante etapas de relés de estado sólido y electromecánicos: paneles de iluminación LED de espectro completo, una red de riego localizado con recolección y recirculación de agua, y ventilación forzada. Piezas a medida diseñadas mediante manufactura aditiva resolvieron los desafíos de impermeabilización, drenaje y montaje mecánico a lo largo de toda la estructura. En cuanto al software, se implementó en Python una lógica de control modular --basada en criterios de umbral e histéresis, filtrado de señal y redundancia de sensores-- para sostener el microclima y el riego dentro de los parámetros objetivo en todo momento, mientras que un algoritmo de presupuesto acumulado regula el fotoperíodo artificial en función de la integral diaria de luz (DLI) recibida por cada cultivo. El prototipo integra además telemetría en la nube y una interfaz web remota, permitiendo el monitoreo continuo y la configuración del sistema conforme al paradigma de Internet de las Cosas (IoT). Los ensayos experimentales realizados a lo largo de varios meses validaron la capacidad del sistema para sostener un microclima estable y favorable para el desarrollo de los cultivos, mientras que las mediciones de consumo de potencia confirmaron que el prototipo opera dentro de un rango energéticamente eficiente, comparable al de una computadora de escritorio estándar.
\vspace{1cm}
\noindent\textbf{Área temática:} Electrónica Digital y Sistemas de Control
\vspace{0.5cm}
\noindent\textbf{Asignaturas:} Electrónica Digital 3, Electrónica Industrial, Electrónica Analógica 1 y Sistemas de Control 1
\vspace{0.5cm}
\noindent\textbf{Palabras clave:} granja vertical, agricultura urbana, riego automatizado, monitoreo ambiental, control embebido, Internet de las Cosas, ciudades inteligentes
